\section{Beispiel}

\subsection{Texte}

Dieser Beispielabschnitt enthält Möglichkeiten, viele wichtige Use-Cases direkt abzuschreiben, anstatt erst Stackoverflow bemühen zu müssen.

Testsätze mit ö, ü und ß.
Zitationen.\cite{Dum:Bei}
Test Email: \email{test@example.com}
Fußnoten.\footnote{Echt cool!}
Eigennamen können so dargestellt werden: \ident{Java}.
Code kann auch mitten im Text dargestellt werden, indem für einfache Sachverhalte \code{monospaced geschrieben wird} und für interessantere Sachen \mint{html}{<h1>mint <b>verwendet<b/> wird</h1>} wodurch dies deutlich spannender aussieht.
Das geht sogar mit Sprachen, die \{ enthalten:
\mint{java}|public static void main (String[] args) { System.out.println("Hello World"); }|

\subsection{Tabellen}

% bessere booktabs-Vergleichstabelle
\begin{table}[H]
    \centering
    \begin{tabular}{>{\raggedleft\arraybackslash}m{0.45\textwidth}>{\raggedright\arraybackslash}m{0.45\textwidth}}
        \toprule
        \centering\arraybackslash\textsc{TESTESTEST} & \centering\arraybackslash\textsc{TESTESTEST} \\\midrule
        \multicolumn{2}{c}{\ident{Beispielheaderzeugs}}                                         \\
        TESTESTESTTESTESTESTTESTES TEST               & TESTESTESTTESTESTES TTESTESTESTTESTESTEST    \\\midrule
        \multicolumn{2}{c}{\ident{Header}}                                                      \\
        asdasdasdasdasdasdasdasdasda                 & a                                            \\\bottomrule
    \end{tabular}
    \caption{Beispiel 2}%
    \label{tab:ex1}
\end{table}

\begin{table}[H]
    \centering
    \begin{tabular}{>{\centering\arraybackslash}m{0.45\textwidth}>{\centering\arraybackslash}m{0.45\textwidth}}
        \toprule
        Stärken                        & Schwächen                      \\\midrule
        TESTESTESTTESTESTESTTESTESTEST & TESTESTESTTESTESTESTTESTESTEST \\
        TESTESTESTTESTESTESTTESTESTEST & TESTESTESTTESTESTESTTESTESTEST \\\bottomrule
    \end{tabular}
    \caption{Beispiel 2}%
    \label{tab:ex2}
\end{table}

% normale Tabelle
\begin{table}[h]
    \centering
    \begin{tabular}{|p{0.25\textwidth}|p{0.25\textwidth}|p{0.25\textwidth}|}
        \hline
        Tabellenkopf 1 & Tabellenkopf 2 & Tabellenkopf 3 \\\hline\hline
        Daten 1        & Daten 2        & Daten 3        \\\hline
        \multicolumn{3}{|c|}{Zusammengefasste Spalten}   \\\hline
    \end{tabular}
    \caption{Eine Beispieltabelle im Text.}
    \label{tab:ex3}
\end{table}

\begin{longtable}{cp{0.6\textwidth}c}
    \caption{Beispiel-LongTable über mehrere Seiten \label{tab:longtable}}\\
    \toprule
    Datum    & Komponente / Aufgabe                                                            & Zeit                    \\* \midrule
    \endfirsthead
    \multicolumn{3}{c}%
    {{\bfseries Tabelle \thetable\ wird von der vorigen Seite fortgeführt.}} \\\midrule
    Datum    & Komponente / Aufgabe                                                            & Zeit                    \\* \midrule
    \endhead
    %
    \bottomrule
    \endfoot
    %
    \bottomrule
    \endlastfoot
    %
    18.10.20 & Organisatorisch (Aufgabenunabhängige Absprache)                                 & 13.30 - 15.00           \\
    19.10.20 & Absprache Aufgabenspezifisch, API Design                                        & 18.00 - 20.30           \\
    20.10.20 & Absprache Entitätenmodellierung Backend                                         & 15.00 - 17.00           \\
    22.10.20 & Absprache + UI Design                                                           & 17.00 - 19.00           \\
    23.10.20 & Initiales SpringBoot-Projekt                                                    & 17.15 - 19.15           \\
    24.10.20 & Pair-Programming JPA/DB-Schicht                                                 & 14.00 - 17.00           \\
    25.10.20 & JPA/DB-Tests                                                                    & 13.30 - 16.30           \\
    25.10.20 & DAO / Spring Data Repositories                                                  & 17.30 - 18.30           \\
    26.10.20 & Besprechung Authentication                                                      & 14.15 - 15.00           \\
    26.10.20 & Diverse Backend Themen                                                          & 16.00 - 18.30           \\
    26.10.20 & Generiere Testdaten                                                             & 18.30 - 19.00           \\
    27.10.20 & Testdaten in Spring integriert                                                  & 09.30 - 10.30           \\
    27.10.20 & ReadMe für das Projekt                                                          & 10.30 - 12.00           \\
    27.10.20 & Review Backend und diverse Absprachen                                           & 17.00 - 17.30           \\
    27.10.20 & Backend-Review Authentication                                                   & 21.30 - 22.15           \\
    28.10.20 & Refactoring Authentication                                                      & 09.30 - 12.00           \\
    28.10.20 & Refactoring Authentication                                                      & 13.00 - 14.30           \\
    29.10.20 & Reviews Docker / Github                                                         & 09.15 - 10.00           \\
    29.10.20 & Backend - Daos                                                                  & 10.00 - 12.00           \\
    29.10.20 & Backend - Daos und Tests                                                        & 13.30 - 17.45           \\
    30.10.20 & Absprache (4-er meeting)                                                        & 14.00 - 15.30
    \\
    02.11.20 & Pair-Programming / Absprache Controller Implementation                          & 11.00 - 15.00           \\
    03.11.20 & Backend - Daos und Tests                                                        & 13.00 - 16.00           \\
    03.11.20 & Lagebesprechung                                                                 & 16.00 - 17.00           \\
    03.11.20 & Backend - Tests für Query und Backend-Besprechung über Error Handling           & 17.00 - 19.00           \\
    05.11.20 & Backend:  Survey-Query, Controller und Service-Schicht + Besprechung und Review & 14.00 - 19.00           \\
    06.11.20 & Backend: Tests für Controller-Schicht                                           & 14.30 - 19.00           \\
    09.11.20 & Backend: Tests für Controller-Schicht                                           & 09.00 - 09.45           \\
    09.11.20 & Code Reviews, Merges und Absprache im Backend-Team für die letzte Woche         & 10.00 - 11.30           \\
    09.11.20 & Backend: Absprachen, Anpassung des Domänenmodells, Refactorings Datenschicht    & 13.00 - 18.00           \\
    09.11.20 & Dokumentation Source Code und Einstellung des Formatters im Backend             & 18.00 - 19.30           \\
    10.11.20 & Dokumentation der Autoren im Backend und Bearbeitung von Linter-Warnungen       & 10.00 - 11.30           \\
    10.11.20 & Schnellabsprache Front\&Backend                                                 & 11.30 - 12.00           \\
    10.11.20 & Arbeit am Abgabedokument                                                        & 15.00 - 19.30           \\
    11.11.20 & Absprache (4-er meeting)                                                        & 16.00 - 17.20 
    \\*
    \end{longtable}

\subsection{Bilder und Minipages}

\begin{figure}[H]
    \centering
    \includegraphics[width=\textwidth]{image/transferleistung.jpg}
    \caption{Ein Beispielbild im Text.}
    \label{fig:tl-logo}
\end{figure}

\begin{figure}[H]
    \begin{minipage}{0.55\textwidth}

        TEST TEST TEST TEST TEST TEST TEST TEST TEST TEST TEST TEST TEST TEST TEST TEST TEST TEST TEST TEST TEST TEST TEST TEST TEST TEST TEST TEST TEST TEST TEST TEST TEST TEST TEST TEST TEST 
        TEST TEST TEST TEST TEST TEST TEST TEST TEST TEST TEST TEST TEST TEST TEST TEST TEST TEST TEST TEST TEST TEST TEST TEST TEST TEST TEST TEST TEST TEST TEST TEST TEST TEST TEST TEST TEST TEST TEST 
        
        TEST TEST TEST TEST TEST TEST TEST TEST TEST TEST TEST TEST TEST TEST TEST TEST TEST TEST TEST TEST TEST TEST TEST TEST TEST TEST TEST TEST TEST TEST TEST TEST TEST TEST TEST TEST TEST TEST TEST TEST TEST TEST TEST TEST TEST TEST TEST TEST TEST TEST TEST TEST TEST TEST TEST TEST TEST TEST TEST 

        TEST TEST TEST TEST TEST TEST TEST TEST TEST TEST TEST TEST TEST TEST TEST TEST TEST TEST TEST TEST TEST TEST TEST TEST TEST TEST TEST TEST TEST TEST TEST TEST TEST TEST TEST TEST TEST TEST TEST TEST TEST TEST TEST TEST TEST TEST TEST TEST TEST TEST TEST TEST TEST TEST TEST TEST TEST TEST TEST TEST TEST TEST TEST TEST TEST TEST TEST TEST TEST TEST 
    \end{minipage}
    \hspace{0.05\textwidth}
    \begin{minipage}{0.4\textwidth}
        \centering
        \includegraphics{image/data-experts-logo.png}
        \caption{Test für Minipages mit korrekter Einrückung ud außerdem mit Bildern}
        \label{fig:ex4}
    \end{minipage}
\end{figure}

\subsection{Modelle}

\begin{figure}[H]
    \centering
    \begin{tikzpicture}[auto,node distance = 0.9cm]
        \node[entity](blatt){AntragsblattImpl}[grow = up, sibling distance = 4cm]
        child {node[key attribute]{id}};

        \node[relationship](rel1)[right = of blatt]{Hat};

        \node[entity](container)[right = of rel1]{TeilVorhabenContainer}[grow = up, sibling distance = 4cm]
        child {node[key attribute]{id}};

        \path(rel1) edge node {1} (blatt) edge node {1} (container);

        \node[relationship](rel2)[right = of container]{Hat};

        \node[entity](content)[right = of rel2]{TeilVorhabenImpl}[grow = up, sibling distance = 2cm]
        child {node[key attribute]{id}}
        child {node[attribute]{owner}};

        \path(rel2) edge node {1} (container) edge node {0-n} (content);
    \end{tikzpicture}
    \caption{Entity-Relationship-Diagramm für das Praxisbeispiel}\label{fig:case:ERD}
\end{figure}

\begin{figure}[H]
    \centering

    \begin{tikzpicture}
        \begin{class}[text width=4cm]{AntragsblattImpl}{0,0}
            \attribute{- id : Long}
            \attribute{- teilVorhabenContainer : TeilVorhabenContainer}
        \end{class}

        \begin{class}[text width=5cm]{TeilVorhabenContainer}{6.5,-0.2}
            \attribute{- id : Long}
            \attribute{- teilvorhaben : List<TeilVorhabenImpl>}
        \end{class}

        \begin{class}[text width=4cm]{TeilVorhabenImpl}{13,-0.65}
            \attribute{- id : Long}
        \end{class}

        \unidirectionalAssociation{AntragsblattImpl}{1}{}{TeilVorhabenContainer}{1}{}
        \unidirectionalAssociation{TeilVorhabenContainer}{1}{}{TeilVorhabenImpl}{0..*}{}
    \end{tikzpicture}

    \caption{UML-Klassendiagramm für das Praxisbeispiel}\label{fig:case:UML}
\end{figure}

\subsection{Codelistings}

\begin{lstlisting}[caption=Sehr simples lstlisting, label=sysout, captionpos=b][H]
    System.out.println("Hallo Welt");
\end{lstlisting}

% works even with π
\begin{listing}[H]
    \begin{minted}[mathescape,
        linenos,
        numbersep=5pt,
        gobble=2,
        frame=lines,
        framesep=2mm]{java}
    String title = "This is a Unicode π in the sky"
    /*
    Defined as $\pi=\lim_{n\to\infty}\frac{P_n}{d}$ where $P$ is the perimeter
    of an $n$-sided regular polygon circumscribing a
    circle of diameter $d$.
    */
    public static void main(String[] args) {
        System.out.println("Hallo, Welt!");
    }
    \end{minted}
    \caption{Ein ausführlicheres Code-Listing mit deutlich mächtigeren Möglichkeiten im Rendering unter der Nutzung von minted}
    \label{code:test}
\end{listing}
